\section*{第 \ref{ch:g12:diversification-of-life} 章——生物的多样化}

\begin{solution}{exo:g12:diversification-of-life:1}
\omterm{def:g10:universal-dna:gene}{基因}复制与分化；水平\omterm{def:g10:universal-dna:gene}{基因}转移；
杂交与多倍化；\omterm{prop:g12:diversification-of-life:development}{发育基因}调控的
改变；\omterm{def:g12:diversification-of-life:symbiosis}{共生}；行为的传递。
\end{solution}

\begin{solution}{exo:g12:diversification-of-life:2}
一组亲缘相关的\omterm{def:g10:universal-dna:gene}{基因}，由同一祖先\omterm{def:g10:universal-dna:gene}{基因}经过连续复制形成，
每个拷贝随后因\omterm{def:g11:mutations:mutation}{突变}而分化，执行不同的功能
（\omterm{def:g11:the-eye:photoreceptors}{视蛋白}、珠蛋白）。
\end{solution}

\begin{solution}{exo:g12:diversification-of-life:3}
其\omterm{def:g11:cell-cycle-mitosis:chromosome}{染色体}来自两个亲本\omterm{def:g10:biodiversity-scales:species}{物种}，每个亲本各提供一套，
在\omterm{def:g12:meiosis-genetic-shuffling:meiosis}{减数分裂}时没有可配对的\omterm{def:g10:common-ancestry:homology}{同源}\omterm{def:g11:cell-cycle-mitosis:chromosome}{染色体}，因此配子不平衡。加倍后，
每条\omterm{def:g11:cell-cycle-mitosis:chromosome}{染色体}都有一条完全相同的伙伴，配对与\omterm{def:g12:meiosis-genetic-shuffling:meiosis}{减数分裂}因而正常进行，
植株可育。
\end{solution}

\begin{solution}{exo:g12:diversification-of-life:4}
拥有自己的环状 \omterm{def:g10:universal-dna:information}{DNA}；\omterm{def:g10:cells-common-unit:organelle}{核糖体}为细菌型，对抑制细菌的\omterm{def:g11:antibiotic-resistance:antibiotic}{抗生素}敏感；
通过分裂增殖；具有双层膜；
其序列与某一类活细菌最为接近。
\end{solution}

\begin{solution}{exo:g12:diversification-of-life:5}
一种由群体中其他成员习得、不依赖\omterm{def:g10:universal-dna:gene}{基因}而代代相传的行为：
某些黑猩猩群体用石头砸坚果，
另一些群体则不然。
\end{solution}

\begin{solution}{exo:g12:diversification-of-life:6}
$14 + 7 = 21$ 条\omterm{def:g11:cell-cycle-mitosis:chromosome}{染色体}：来自一粒小麦的 7 条 A \omterm{def:g11:cell-cycle-mitosis:chromosome}{染色体}，以及来自二粒小麦的 7 条 A 和 7 条 B \omterm{def:g11:cell-cycle-mitosis:chromosome}{染色体}。
7 对（A 与 A）可以配成对；7 条 B \omterm{def:g11:cell-cycle-mitosis:chromosome}{染色体}没有配对的伙伴，
在分配时随机分开，因此配子不平衡：不育。
\end{solution}

\begin{solution}{exo:g12:diversification-of-life:7}
$\alpha$--$\beta$ 的分歧（序列 50\% 相同）早于
$\beta$--$\gamma$ 的分歧（80\%）：$\alpha$ 最先分支，然后 $\beta$
与 $\gamma$ 相互分开。
\end{solution}

\begin{solution}{exo:g12:diversification-of-life:8}
肢体\omterm{def:g10:universal-dna:gene}{基因}本身是存在的，但在胚胎侧翼开启它们的信号
未曾出现，或者出现后又中断；
两者的差异在于\omterm{def:g10:universal-dna:gene}{基因}表达的调控，而不在\omterm{def:g10:universal-dna:gene}{基因}本身。
\end{solution}

\begin{solution}{exo:g12:diversification-of-life:9}
差异在于胚胎喙部生长信号的强度与出现时间：
这是一种调控上的差异。调控的改变只需
控制序列中发生一个小\omterm{def:g11:mutations:mutation}{突变}，而一种新\omterm{def:g11:gene-expression:protein}{蛋白质}则需要许多\omterm{def:g11:mutations:mutation}{突变}；
它可以在数千代之内产生并被选择。
\end{solution}

\begin{solution}{exo:g12:diversification-of-life:10}
其形态、在裸露岩石上定居的能力、对干旱的抗性，
以及缓慢的生长，都不属于任何一方单独拥有，
而是来自二者的联合。地衣是一种复合生物体，
具有其自身的特性。
\end{solution}

\begin{solution}{exo:g12:diversification-of-life:11}
行为的传递：两岸都有相同的材料，
种群之间亲缘关系密切，差异恰好以
河流为界——河流是接触的障碍，因而也是学习的障碍，
而不是\omterm{def:g10:universal-dna:gene}{基因}或资源的障碍。
\end{solution}

\begin{solution}{exo:g12:diversification-of-life:12}
一种病毒在数千万年前将其\omterm{def:g10:universal-dna:gene}{基因}插入灵长类祖先的
生殖\omterm{def:g10:cells-common-unit:cell}{细胞}的\omterm{def:g10:universal-dna:gene}{基因}组中；该\omterm{def:g10:universal-dna:gene}{基因}被遗传下来并保留至今。
检验方法：该\omterm{def:g10:universal-dna:gene}{基因}应位于所有灵长类的同一\omterm{def:g11:cell-cycle-mitosis:chromosome}{染色体}位置，
并在其他哺乳动物的该位置上缺失；
其序列应与该病毒的序列最为接近。
\end{solution}

\begin{solution}{exo:g12:diversification-of-life:13}
\omterm{def:g10:cells-common-unit:organelle}{线粒体}\omterm{def:g10:cells-common-unit:organelle}{核糖体}属于细菌型，并与同样的\omterm{def:g11:antibiotic-resistance:antibiotic}{抗生素}结合：
高剂量时，该药物会减缓\omterm{def:g10:cells-common-unit:organelle}{线粒体}的\omterm{def:g11:gene-expression:protein}{蛋白质}合成。
这种敏感性遗传自\omterm{def:g10:cells-common-unit:organelle}{线粒体}的
细菌祖先。
\end{solution}

\begin{solution}{exo:g12:diversification-of-life:14}
\omterm{prop:g12:diversification-of-life:polyploidy}{多倍体}的\omterm{def:g11:cell-cycle-mitosis:chromosome}{染色体}数目与任一亲本都不相同，
因此它们与亲本杂交的后代不育：生殖隔离即刻形成。
一个新的\omterm{def:g10:universal-dna:gene}{等位基因}需要经过许多代才能在种群中扩散，
才能使某一类群得以区分。因此，一个植物\omterm{def:g10:biodiversity-scales:species}{物种}可以
在一代之内产生。
\end{solution}

\begin{solution}{exo:g12:diversification-of-life:15}
归根结底，每一段新序列——无论是复制后分化的拷贝，
调控上的改变，还是转移质粒上的\omterm{def:g10:universal-dna:gene}{基因}——最初都源于某处的一次\omterm{def:g11:mutations:mutation}{突变}。
但是，复制、转移与多倍化能同时迁移
并增殖整个\omterm{def:g10:universal-dna:gene}{基因}与\omterm{def:g10:universal-dna:gene}{基因}组；调控的改变能够产生
无需新\omterm{def:g11:gene-expression:protein}{蛋白质}的新形态；\omterm{def:g12:diversification-of-life:symbiosis}{共生}则将
不同\omterm{def:g10:biodiversity-scales:species}{物种}的\omterm{def:g10:universal-dna:gene}{基因}组结合在一起；文化则在完全不依赖\omterm{def:g10:universal-dna:gene}{基因}的
情况下传递变异。\omterm{def:g11:mutations:mutation}{突变}提供了字母，这些过程则重新排列
整页乃至整本书。
\end{solution}

\begin{solution}{pb:g12:diversification-of-life:1}
\textbf{1.} 二粒小麦 28，面包小麦 42。

\textbf{2.} 没有：A 与 B \omterm{def:g11:cell-cycle-mitosis:chromosome}{染色体}不是\omterm{def:g10:common-ancestry:homology}{同源}\omterm{def:g11:cell-cycle-mitosis:chromosome}{染色体}。配子
只能随机获得 14 条中的一部分，几乎不可能得到完整的一套。

\textbf{3.} 共 28 条\omterm{def:g11:cell-cycle-mitosis:chromosome}{染色体}，每条 A 都与一条相同的 A 配对，
每条 B 都与一条 B 配对：14 对，\omterm{def:g12:meiosis-genetic-shuffling:meiosis}{减数分裂}正常，配子各含 14 条。

\textbf{4.} 均不能配对：7 条 A、7 条 B 和 7 条 D \omterm{def:g11:cell-cycle-mitosis:chromosome}{染色体}，每条都
缺少\omterm{def:g10:common-ancestry:homology}{同源}\omterm{def:g11:cell-cycle-mitosis:chromosome}{染色体}，不育。

\textbf{5.} 21 对；花粉粒携带 21 条。

\textbf{6.} 另外两个拷贝仍能提供该\omterm{def:g11:gene-expression:protein}{蛋白质}。冗余
拷贝可以自由积累\omterm{def:g11:mutations:mutation}{突变}并发生分化——重复的范围
达到整个\omterm{def:g10:universal-dna:gene}{基因}组的规模。

\textbf{7.} 在二粒小麦种植地区的田间，
距今 10\,000 至 8\,000 年，野生山羊草 2 作为杂草
混生于作物之中。

\textbf{8.} 说明两个步骤——杂交与加倍——足以
由亲本产生面包小麦：这段历史可以
重新演示。

\textbf{9.} 是：它与单粒小麦无法产生可育后代，所以
依定义它是一个独立\omterm{def:g10:biodiversity-scales:species}{物种}，尽管它是
由单粒小麦演化而来的。

\textbf{10.} 野生山羊草 2 的 D \omterm{def:g11:cell-cycle-mitosis:chromosome}{染色体}与小麦的 D \omterm{def:g11:cell-cycle-mitosis:chromosome}{染色体}组\omterm{def:g10:common-ancestry:homology}{同源}，
能与之配对，因此杂交并经回交即可将该\omterm{def:g10:universal-dna:gene}{基因}导入；
无亲缘关系植物的\omterm{def:g11:cell-cycle-mitosis:chromosome}{染色体}则根本无法配对。

\textbf{11.} \omterm{def:g10:universal-dna:gene}{基因}本身几乎没有差别；身体的差异
必定在于\omterm{def:g10:universal-dna:gene}{基因}的使用方式。

\textbf{12.} 表达调控的改变——即
\omterm{def:g10:universal-dna:gene}{基因}在体轴上活跃区域的范围。

\textbf{13.} 四肢\omterm{def:g10:universal-dna:gene}{基因}存在且开始发挥作用；维持四肢生长的
信号缺失，因此芽停止生长。同样，这是调控的问题，
而非\omterm{def:g10:universal-dna:gene}{基因}的丢失。

\textbf{14.} 这表明，改变\omterm{def:g10:universal-dna:gene}{基因}的表达位置
就足以改变身体区域：表达与形态之间的相关性
是因果关系。

\textbf{15.} 面包小麦的变化是为\omterm{def:g10:universal-dna:gene}{基因}组增添了整套\omterm{def:g11:cell-cycle-mitosis:chromosome}{染色体}（即
\omterm{def:g10:universal-dna:gene}{基因}组大小）；蟒蛇的变化则在于共有\omterm{def:g10:universal-dna:gene}{基因}被使用的位置（即
其调控方式）。

\textbf{16.} 自身的环状 \omterm{def:g10:universal-dna:information}{DNA}（预期：无，全部\omterm{def:g10:universal-dna:gene}{基因}位于\omterm{def:g10:cells-common-unit:organelle}{细胞核}中）；
细菌型\omterm{def:g10:cells-common-unit:organelle}{核糖体}（预期：真\omterm{def:g11:cell-cycle-mitosis:chromosome}{核型}\omterm{def:g10:cells-common-unit:organelle}{核糖体}）；通过
分裂增殖（预期：由组件装配）；双层膜
（预期：与其他\omterm{def:g10:cells-common-unit:organelle}{细胞器}一样为单层膜）。而且它们的\omterm{def:g10:universal-dna:gene}{基因}
类似细菌\omterm{def:g10:universal-dna:gene}{基因}（预期：\omterm{def:g10:cells-common-unit:organelle}{细胞核}型\omterm{def:g10:universal-dna:gene}{基因}）。

\textbf{17.} 在\omterm{def:g10:cells-common-unit:organelle}{细胞核}中：随着时间推移，细菌的大部分\omterm{def:g10:universal-dna:gene}{基因}
已转移至\omterm{def:g10:cells-common-unit:organelle}{细胞核}\omterm{def:g10:universal-dna:gene}{基因}组，而该\omterm{def:g10:cells-common-unit:organelle}{细胞器}自身的\omterm{def:g10:universal-dna:gene}{基因}组
缩减为残余。

\textbf{18.} \omterm{def:g10:cells-common-unit:organelle}{线粒体}\omterm{def:g10:universal-dna:gene}{基因}只由母亲传给子女；一个人的\omterm{def:g10:cells-common-unit:organelle}{线粒体}
\omterm{def:g10:universal-dna:information}{DNA} 追溯的是未中断的母系谱系。

\textbf{19.} \omterm{def:g10:cells-common-unit:organelle}{线粒体}已将大部分\omterm{def:g10:universal-dna:gene}{基因}转交给宿主，无法
脱离\omterm{def:g10:cells-common-unit:cell}{细胞}而存在；地衣的\omterm{def:g12:diversification-of-life:symbiosis}{共生}伙伴则保留了自身\omterm{def:g10:universal-dna:gene}{基因}组，
并可在一定困难下被分离。

\textbf{20.} 单粒小麦（14 条，A）、野生山羊草 1（14 条，B）、
野生山羊草 2（14 条，D）：共 42 条\omterm{def:g11:cell-cycle-mitosis:chromosome}{染色体}；“调控”；\omterm{def:g10:cells-common-unit:organelle}{线粒体}。
\end{solution}
